% Slot 3 (Fri 27 Nov 10:15--12:00 Rome)

\section{Slot 3}

\begin{frame}{Slot 3 \textperiodcentered\ Fri 10:15--12:00}
  \begin{center}
    \begin{tabular}{@{}rl@{}}
      \toprule
      \textbf{Clock} & \textbf{Block} \\
      \midrule
      10:15--10:50 & Using sub-agents \\
      10:50--11:15 & Orchestrating sub-agents and swarms \\
      11:15--11:25 & Pause \\
      11:25--12:00 & Issue-board swarm: launch two cards \\
      \textit{12:00--14:00} & \textit{Lunch: the agents keep working} \\
      \bottomrule
    \end{tabular}
  \end{center}
\end{frame}

% ----------------------------------------------------------------
\begin{frame}{Using sub-agents}
  \timebox{10:15--10:50}
  \inplainwords{Name the sub-agent and the bounds for this piece of work. It works
    in its own conversation and hands back what the job needs. These two reviewers
    may only read; a review is what they hand back. You decide what to keep.}
  \likethis{two research assistants. One reviews the cleaning, one checks the
    citations. Neither sees the other's notes.}
\end{frame}

\begin{frame}{Card--Krueger: two reviews}
  \prompttext{\raggedright Use the data-reviewer sub-agent on examples/card-krueger-toy/. Read only. Report panel balance and NJ/PA coding.}
  \vspace{0.4em}
  \prompttext{\raggedright Use the literature-reviewer sub-agent on the same folder. Read only. Check the claims against the references.}
  \vspace{0.6em}
  \begin{itemize}
    \item Two separate conversations: neither sees the other's review.
    \item Where they agree, look closer. Where they disagree, you decide.
  \end{itemize}
\end{frame}

\begin{frame}{Your turn}
  \youtry{10:15--10:50}
  \vspace{0.3em}
  \begin{enumerate}
    \item Run both reviewers on the Card--Krueger folder.
    \item Log one finding you accept and one you reject in \artifact{notes/orchestration\_log.md}.
    \item Check that no file changed.
  \end{enumerate}
  \vspace{0.4em}
  {\small Reviewer won't start? Paste its instructions into a fresh chat and say ``read only''.}\\
  {\small Guide: \artifact{demos/slot3\_subagents\_vc.md} \textperiodcentered\ \artifact{exercises/07\_slot3\_subagents\_vc.md}}
\end{frame}

% ----------------------------------------------------------------
\begin{frame}{Orchestrating sub-agents and swarms}
  \timebox{10:50--11:15}
  \inplainwords{Orchestrating means splitting a job into parts, handing each part to a
    different agent, and putting the results back together. A swarm is many agents
    working at the same time; a human still makes the final call.}
  \likethis{a PI running several RAs, each with their own files.}
\end{frame}

\begin{frame}{Card--Krueger: a writer and a checker}
  \begin{center}
    \begin{tikzpicture}[
        node distance=1.0cm,
        role/.style={rectangle, draw=courseblue, thick, rounded corners=3pt, minimum width=3.3cm, minimum height=1.0cm, align=center, font=\footnotesize\bfseries, fill=coursebg},
        arr/.style={-Stealth, thick, courseblue}
      ]
      \node[role] (w) {Writer\\\normalfont\scriptsize documents the panel};
      \node[role, right=of w, fill=courseteal!14] (c) {Checker\\\normalfont\scriptsize reviews literature claims};
      \node[role, below=0.9cm of $(w)!0.5!(c)$, draw=codered, fill=codered!8] (y) {You\\\normalfont\scriptsize read both, then merge};
      \draw[arr] (w) -- (y);
      \draw[arr] (c) -- (y);
    \end{tikzpicture}
  \end{center}
  \vspace{0.3em}
  \begin{itemize}
    \item Different files: the two agents never edit the same lines.
    \item Rules for the pair: \artifact{examples/card-krueger-swarm/README.md}.
  \end{itemize}
\end{frame}

\begin{frame}{Your turn: split the job}
  \youtry{10:50--11:15}
  \vspace{0.3em}
  Split ``replicate and extend Card--Krueger'' into three jobs that share no files.
  \vspace{0.5em}

  \begin{tabular}{@{}lp{0.40\linewidth}p{0.32\linewidth}@{}}
    \toprule
    \textbf{Job} & \textbf{Files it may touch} & \textbf{Done when} \\
    \midrule
    1 & & \\
    2 & & \\
    3 & & \\
    \bottomrule
  \end{tabular}
  \vspace{0.5em}

  {\small Pause 11:15--11:25.}
\end{frame}

% ----------------------------------------------------------------
\begin{frame}{Issue-board swarm}
  \timebox{11:25--12:00}
  \inplainwords{Each job is written as an issue, a to-do card on GitHub. Agents pick
    up cards marked ready, work in their own copy, and send back a pull request
    linked to the card for review.}
  \likethis{a job board in a shared office: whoever is free takes the next card.}
\end{frame}

\begin{frame}{Card--Krueger: two cards on the board}
  \begin{center}
    \begin{tikzpicture}[
        node distance=0.55cm,
        st/.style={rectangle, draw=courseblue, thick, rounded corners=3pt, minimum width=2.4cm, minimum height=0.8cm, align=center, font=\footnotesize\bfseries, fill=coursebg},
        arr/.style={-Stealth, thick, courseblue}
      ]
      \node[st] (r) {ready};
      \node[st, right=of r] (p) {in-progress};
      \node[st, right=of p, fill=courseteal!14] (n) {needs-review};
      \node[st, right=of n, fill=codegreen!12] (d) {done};
      \draw[arr] (r) -- (p);
      \draw[arr] (p) -- (n);
      \draw[arr] (n) -- (d);
    \end{tikzpicture}
  \end{center}
  \vspace{0.4em}
  \begin{itemize}
    \item Card 1: \artifact{ck-data-cleaning}: check and document the panel.
    \item Card 2: \artifact{ck-parallel-trends-figure}: draw the before/after figure.
    \item Cards live in \artifact{examples/card-krueger-swarm/issues/}.
    \item Only you move a card to \textbf{done}.
  \end{itemize}
\end{frame}

\begin{frame}{Your turn: launch before lunch}
  \youtry{11:25--12:00}
  \vspace{0.3em}
  \begin{enumerate}
    \item Cloud agent? First fill in \artifact{agent-harness/autonomous\_agent\_risk\_card.md}.
    \item Open the two cards as issues and mark them \textbf{ready}.
    \item Start two agents, one per card (cloud, or two local windows).
    \item Go to lunch. Back at 14:00 to review.
  \end{enumerate}
  \vspace{0.4em}
  {\small At most two cards live. No confidential files. Agents never merge.}
\end{frame}
